\section{A Note for Readers Who Struggle with Sleep}\label{sec:note_for_reader}

\textbf{If you struggle with sleep \- or even dread going to bed \- this report is not here to scold you.}

We just reviewed stark findings on sleep loss. Those are real. But most readers don`t lack knowledge; they lack \emph{access} to sleep. In real life, nights can feel unsafe, lonely, or unbearably boring; ``just sleep more’’ can land like blame rather than help.

This isn`t laziness or a moral failing. It`s a system under strain \citep{sateia2014icsd3_overview}.

\textbf{Why sleep can be hard (non-exhaustive):}
\begin{itemize}
  \item \textbf{Mental health:} Anxiety, depression, and PTSD commonly co-occur with insomnia and nightmares; hyperarousal itself is part of PTSD`s diagnosis \citep{lancel2021ptsd_sleep,morgenthaler2018aasm_nightmares}.
  \item \textbf{Neurodivergence:} ADHD is frequently linked to delayed circadian phase and sleep-onset insomnia; stimulants can further delay circadian timing or cause insomnia as a side effect \citep{snitselaar2017adhd_circadian,auger2015aasm_circadian}.
  \item \textbf{Substance use (including prescribed stimulants):} Chronic stimulant exposure is associated with clinically significant insomnia that often needs active management \citep{chiappini2025stimulant_insomnia_mgmt,oliva2025stimulants_safety_meta}.
  \item \textbf{Chronic pain and medical conditions:} Pain, OSA, and RLS are tightly bidirectionally linked with sleep disturbance \citep{duan2024pain_sleep_bidirectional}.
  \item \textbf{Circadian rhythm disorders:} In Delayed Sleep\-Wake Phase Disorder, internal time is shifted later; forcing “normal” hours can feel like perpetual jet lag \citep{auger2015aasm_circadian}.
  \item \textbf{Trauma and hypervigilance:} The nervous system may treat sleep as loss of control, keeping you on guard \citep{sateia2014icsd3_overview}.
\end{itemize}

If you live with ADHD, trauma, circadian delay, chronic pain, or you use stimulants, struggling with sleep is not a character flaw or “poor sleep hygiene.” Major guidelines are clear: for chronic insomnia, tips-and-tricks alone are \textbf{not} an effective treatment; it`s a medical/behavioral condition best addressed with structured care (e.g., CBT-I), and circadian disorders (like delayed sleep phase) often need targeted timing, light, and/or melatonin strategies. Stimulants can biologically delay sleep or cause insomnia as a side effect. In other words: you are not lazy; your system is under specific, workable constraints, and there are evidence-based ways to help \citep{qaseem2016acp_insomnia,schutte2008aasm_insomnia,auger2015aasm_circadian,snitselaar2013adhd_circadian,fda2023vyvanse_label}.

